%% rset.future.tex
%% Future work: arbitrarily many colours, the finite-state-machine
%% generalization, and the open problem of generating colours dynamically.
%% Drop in via %% rset.future.tex
%% Future work: arbitrarily many colours, the finite-state-machine
%% generalization, and the open problem of generating colours dynamically.
%% Drop in via %% rset.future.tex
%% Future work: arbitrarily many colours, the finite-state-machine
%% generalization, and the open problem of generating colours dynamically.
%% Drop in via %% rset.future.tex
%% Future work: arbitrarily many colours, the finite-state-machine
%% generalization, and the open problem of generating colours dynamically.
%% Drop in via \input{rset.future}.
\input{rset.colors}   % red/black colour layer + FM symbols (providecommand-guarded)

\section{Future work: palettes and the generation of colour}
\label{sec:future}

\subsection{Arbitrarily many colours}
\label{sec:future:palette}

Two colours were a convenience, not a necessity. Replace the pair by any set $C$
of \emph{colours}; each $c\in C$ carries its own copy of the theory,
$\mathrm{Set}_c[X_c]$, with its own membership $\in_c$ and its own comprehension
$\{\cdots\}_c$. The only remaining datum is the \emph{nesting discipline}: which
colours may serve as atoms for which. This is a relation on $C$, and it is most
naturally presented as a directed graph --- equivalently a finite-state machine
--- with states $C$ and an edge $c\to c'$ whenever $c'$-sets are admissible
$c$-atoms. The atom supply of each colour is then the disjoint sum over its
outgoing edges,
\[
  X_c \;=\; \bigsqcup_{c\to c'}\ \mathrm{Set}_{c'}[X_{c'}],
\]
and the two-colour theory of this paper is the instance given by the two-state
automaton $\rd{\mathrm{red}} \rightleftarrows \bk{\mathrm{black}}$, with
$\rX=\bSetX$ and $\bX=\rSetX$.

Everything in the paper relativises to this picture. The well-formed cross-colour
nestings are the runs of the automaton; the local/global well-foundedness analysis
of Section~\ref{sec:nwf} reads off its structure, an acyclic colour graph giving a
globally well-founded universe and a cycle giving non-well-foundedness along that
cycle. (A self-loop $c\to c$ would let $c$-sets be their own atoms and collapse
the barrier; the two-cycle $\rd{\mathrm{red}}\rightleftarrows\bk{\mathrm{black}}$
is the smallest discipline that keeps every colour locally well-founded while
making the system globally non-well-founded.) This generalization is routine, and
we record it mainly to fix the slogan: \emph{a reflective set theory is an
automaton on colours}, and the constructions above are its two-state instance.

\subsection{Generating colours dynamically}
\label{sec:future:dynamic}

The genuinely interesting question we leave open. A fixed automaton fixes the
palette in advance. One would instead like colours to be \emph{generated on
demand}: a fresh colour spawned in the course of building a set, the palette
growing as construction proceeds, the automaton unfolding rather than given. The
difficulty is not in admitting infinitely many colours --- a countable palette is
unremarkable --- but in doing so \emph{elegantly}, with a generation discipline
internal to the theory, so that the mechanism producing colours is of a piece with
the mechanism producing sets and names.

There is a suggestive circularity to exploit, which is why we think the question
worth posing rather than dismissing. Generating a fresh colour is generating a
fresh name, in exactly the sense that
Sections~\ref{sec:fm}--\ref{sec:applications} reconstruct: a colour is drawn from
a supply, distinct from those already in play, and otherwise structureless. So a
reflective theory might be made to \emph{generate its own colours} by the very
nominal machinery it models --- the palette drawn, like the atoms, from a stream,
and the colour automaton itself becoming a nominal object in a theory one level
up. Whether this reflexive step can be taken without infinite regress, what its
fixed point looks like, and whether it lands back inside the spectrum of
Section~\ref{sec:nwf}, are the questions we find most worth pursuing. We do not
settle them here; we mark them as the natural sequel to this work.
.
%% rset.colors.tex
%% Red/black colour layer for the two intertwined copies, plus the FM symbols
%% used by rset.mainthm.tex and rset.applications.tex.
%%
%% EVERYTHING routes through \rd (red copy = player) and \bk (black copy =
%% opponent): redefine just those two macros to change the whole scheme (e.g.
%% to make the "black" copy dark blue for legibility).  Every definition is
%% \providecommand-guarded, so it is safe to \input this file from several
%% places and harmless if rset.local already defines these names.
%%
%% Requires \usepackage{xcolor} (already loaded in rset.tex).

% --- the two copies -------------------------------------------------------
\providecommand{\rd}[1]{\textcolor{red}{#1}}     % red copy  (player)
\providecommand{\bk}[1]{\textcolor{black}{#1}}   % black copy (opponent)

% --- theory names and atom supplies --------------------------------------
\providecommand{\rSetX}{\rd{\mathrm{Set}[X]}}    % red   Set[X]
\providecommand{\bSetX}{\bk{\mathrm{Set}[X]}}    % black Set[X]
\providecommand{\rX}{\rd{X}}                      % red   atom supply
\providecommand{\bX}{\bk{X}}                      % black atom supply

% --- empty set ------------------------------------------------------------
\providecommand{\remp}{\rd{\emptyset}}
\providecommand{\bemp}{\bk{\emptyset}}

% --- membership (relations) ----------------------------------------------
\providecommand{\rin}{\mathrel{\rd{\in}}}
\providecommand{\bin}{\mathrel{\bk{\in}}}
\providecommand{\rnin}{\mathrel{\rd{\notin}}}
\providecommand{\bnin}{\mathrel{\bk{\notin}}}

% --- binary operations ----------------------------------------------------
\providecommand{\rcup}{\mathbin{\rd{\cup}}}
\providecommand{\bcup}{\mathbin{\bk{\cup}}}
\providecommand{\rcap}{\mathbin{\rd{\cap}}}
\providecommand{\bcap}{\mathbin{\bk{\cap}}}
\providecommand{\rsm}{\mathbin{\rd{\setminus}}}
\providecommand{\bsm}{\mathbin{\bk{\setminus}}}

% --- comprehension braces of each theory ---------------------------------
\providecommand{\rbr}[1]{\rd{\{}\,#1\,\rd{\}}}
\providecommand{\bbr}[1]{\bk{\{}\,#1\,\bk{\}}}

% --- colour-blind membership (sees through the colour barrier) ------------
\providecommand{\cbin}{\mathrel{\in_{\!\star}}}

% --- FM symbols (shared with rset.stdfmset.tex) --------------------------
\providecommand{\Atoms}{\mathbb{A}}
\providecommand{\Sym}{\mathrm{Sym}}
\providecommand{\pow}{\mathcal{P}}
\providecommand{\supp}{\mathrm{supp}}
\providecommand{\fs}{\mathrm{fs}}
\providecommand{\new}{\mathbin{\reflectbox{\ensuremath{\mathsf{N}}}}}
\providecommand{\atomsof}{\mathrm{atoms}}
\providecommand{\Red}{\rd{\mathsf{R}}}           % the red HF universe
   % red/black colour layer + FM symbols (providecommand-guarded)

\section{Future work: palettes and the generation of colour}
\label{sec:future}

\subsection{Arbitrarily many colours}
\label{sec:future:palette}

Two colours were a convenience, not a necessity. Replace the pair by any set $C$
of \emph{colours}; each $c\in C$ carries its own copy of the theory,
$\mathrm{Set}_c[X_c]$, with its own membership $\in_c$ and its own comprehension
$\{\cdots\}_c$. The only remaining datum is the \emph{nesting discipline}: which
colours may serve as atoms for which. This is a relation on $C$, and it is most
naturally presented as a directed graph --- equivalently a finite-state machine
--- with states $C$ and an edge $c\to c'$ whenever $c'$-sets are admissible
$c$-atoms. The atom supply of each colour is then the disjoint sum over its
outgoing edges,
\[
  X_c \;=\; \bigsqcup_{c\to c'}\ \mathrm{Set}_{c'}[X_{c'}],
\]
and the two-colour theory of this paper is the instance given by the two-state
automaton $\rd{\mathrm{red}} \rightleftarrows \bk{\mathrm{black}}$, with
$\rX=\bSetX$ and $\bX=\rSetX$.

Everything in the paper relativises to this picture. The well-formed cross-colour
nestings are the runs of the automaton; the local/global well-foundedness analysis
of Section~\ref{sec:nwf} reads off its structure, an acyclic colour graph giving a
globally well-founded universe and a cycle giving non-well-foundedness along that
cycle. (A self-loop $c\to c$ would let $c$-sets be their own atoms and collapse
the barrier; the two-cycle $\rd{\mathrm{red}}\rightleftarrows\bk{\mathrm{black}}$
is the smallest discipline that keeps every colour locally well-founded while
making the system globally non-well-founded.) This generalization is routine, and
we record it mainly to fix the slogan: \emph{a reflective set theory is an
automaton on colours}, and the constructions above are its two-state instance.

\subsection{Generating colours dynamically}
\label{sec:future:dynamic}

The genuinely interesting question we leave open. A fixed automaton fixes the
palette in advance. One would instead like colours to be \emph{generated on
demand}: a fresh colour spawned in the course of building a set, the palette
growing as construction proceeds, the automaton unfolding rather than given. The
difficulty is not in admitting infinitely many colours --- a countable palette is
unremarkable --- but in doing so \emph{elegantly}, with a generation discipline
internal to the theory, so that the mechanism producing colours is of a piece with
the mechanism producing sets and names.

There is a suggestive circularity to exploit, which is why we think the question
worth posing rather than dismissing. Generating a fresh colour is generating a
fresh name, in exactly the sense that
Sections~\ref{sec:fm}--\ref{sec:applications} reconstruct: a colour is drawn from
a supply, distinct from those already in play, and otherwise structureless. So a
reflective theory might be made to \emph{generate its own colours} by the very
nominal machinery it models --- the palette drawn, like the atoms, from a stream,
and the colour automaton itself becoming a nominal object in a theory one level
up. Whether this reflexive step can be taken without infinite regress, what its
fixed point looks like, and whether it lands back inside the spectrum of
Section~\ref{sec:nwf}, are the questions we find most worth pursuing. We do not
settle them here; we mark them as the natural sequel to this work.
.
%% rset.colors.tex
%% Red/black colour layer for the two intertwined copies, plus the FM symbols
%% used by rset.mainthm.tex and rset.applications.tex.
%%
%% EVERYTHING routes through \rd (red copy = player) and \bk (black copy =
%% opponent): redefine just those two macros to change the whole scheme (e.g.
%% to make the "black" copy dark blue for legibility).  Every definition is
%% \providecommand-guarded, so it is safe to \input this file from several
%% places and harmless if rset.local already defines these names.
%%
%% Requires \usepackage{xcolor} (already loaded in rset.tex).

% --- the two copies -------------------------------------------------------
\providecommand{\rd}[1]{\textcolor{red}{#1}}     % red copy  (player)
\providecommand{\bk}[1]{\textcolor{black}{#1}}   % black copy (opponent)

% --- theory names and atom supplies --------------------------------------
\providecommand{\rSetX}{\rd{\mathrm{Set}[X]}}    % red   Set[X]
\providecommand{\bSetX}{\bk{\mathrm{Set}[X]}}    % black Set[X]
\providecommand{\rX}{\rd{X}}                      % red   atom supply
\providecommand{\bX}{\bk{X}}                      % black atom supply

% --- empty set ------------------------------------------------------------
\providecommand{\remp}{\rd{\emptyset}}
\providecommand{\bemp}{\bk{\emptyset}}

% --- membership (relations) ----------------------------------------------
\providecommand{\rin}{\mathrel{\rd{\in}}}
\providecommand{\bin}{\mathrel{\bk{\in}}}
\providecommand{\rnin}{\mathrel{\rd{\notin}}}
\providecommand{\bnin}{\mathrel{\bk{\notin}}}

% --- binary operations ----------------------------------------------------
\providecommand{\rcup}{\mathbin{\rd{\cup}}}
\providecommand{\bcup}{\mathbin{\bk{\cup}}}
\providecommand{\rcap}{\mathbin{\rd{\cap}}}
\providecommand{\bcap}{\mathbin{\bk{\cap}}}
\providecommand{\rsm}{\mathbin{\rd{\setminus}}}
\providecommand{\bsm}{\mathbin{\bk{\setminus}}}

% --- comprehension braces of each theory ---------------------------------
\providecommand{\rbr}[1]{\rd{\{}\,#1\,\rd{\}}}
\providecommand{\bbr}[1]{\bk{\{}\,#1\,\bk{\}}}

% --- colour-blind membership (sees through the colour barrier) ------------
\providecommand{\cbin}{\mathrel{\in_{\!\star}}}

% --- FM symbols (shared with rset.stdfmset.tex) --------------------------
\providecommand{\Atoms}{\mathbb{A}}
\providecommand{\Sym}{\mathrm{Sym}}
\providecommand{\pow}{\mathcal{P}}
\providecommand{\supp}{\mathrm{supp}}
\providecommand{\fs}{\mathrm{fs}}
\providecommand{\new}{\mathbin{\reflectbox{\ensuremath{\mathsf{N}}}}}
\providecommand{\atomsof}{\mathrm{atoms}}
\providecommand{\Red}{\rd{\mathsf{R}}}           % the red HF universe
   % red/black colour layer + FM symbols (providecommand-guarded)

\section{Future work: palettes and the generation of colour}
\label{sec:future}

\subsection{Arbitrarily many colours}
\label{sec:future:palette}

Two colours were a convenience, not a necessity. Replace the pair by any set $C$
of \emph{colours}; each $c\in C$ carries its own copy of the theory,
$\mathrm{Set}_c[X_c]$, with its own membership $\in_c$ and its own comprehension
$\{\cdots\}_c$. The only remaining datum is the \emph{nesting discipline}: which
colours may serve as atoms for which. This is a relation on $C$, and it is most
naturally presented as a directed graph --- equivalently a finite-state machine
--- with states $C$ and an edge $c\to c'$ whenever $c'$-sets are admissible
$c$-atoms. The atom supply of each colour is then the disjoint sum over its
outgoing edges,
\[
  X_c \;=\; \bigsqcup_{c\to c'}\ \mathrm{Set}_{c'}[X_{c'}],
\]
and the two-colour theory of this paper is the instance given by the two-state
automaton $\rd{\mathrm{red}} \rightleftarrows \bk{\mathrm{black}}$, with
$\rX=\bSetX$ and $\bX=\rSetX$.

Everything in the paper relativises to this picture. The well-formed cross-colour
nestings are the runs of the automaton; the local/global well-foundedness analysis
of Section~\ref{sec:nwf} reads off its structure, an acyclic colour graph giving a
globally well-founded universe and a cycle giving non-well-foundedness along that
cycle. (A self-loop $c\to c$ would let $c$-sets be their own atoms and collapse
the barrier; the two-cycle $\rd{\mathrm{red}}\rightleftarrows\bk{\mathrm{black}}$
is the smallest discipline that keeps every colour locally well-founded while
making the system globally non-well-founded.) This generalization is routine, and
we record it mainly to fix the slogan: \emph{a reflective set theory is an
automaton on colours}, and the constructions above are its two-state instance.

\subsection{Generating colours dynamically}
\label{sec:future:dynamic}

The genuinely interesting question we leave open. A fixed automaton fixes the
palette in advance. One would instead like colours to be \emph{generated on
demand}: a fresh colour spawned in the course of building a set, the palette
growing as construction proceeds, the automaton unfolding rather than given. The
difficulty is not in admitting infinitely many colours --- a countable palette is
unremarkable --- but in doing so \emph{elegantly}, with a generation discipline
internal to the theory, so that the mechanism producing colours is of a piece with
the mechanism producing sets and names.

There is a suggestive circularity to exploit, which is why we think the question
worth posing rather than dismissing. Generating a fresh colour is generating a
fresh name, in exactly the sense that
Sections~\ref{sec:fm}--\ref{sec:applications} reconstruct: a colour is drawn from
a supply, distinct from those already in play, and otherwise structureless. So a
reflective theory might be made to \emph{generate its own colours} by the very
nominal machinery it models --- the palette drawn, like the atoms, from a stream,
and the colour automaton itself becoming a nominal object in a theory one level
up. Whether this reflexive step can be taken without infinite regress, what its
fixed point looks like, and whether it lands back inside the spectrum of
Section~\ref{sec:nwf}, are the questions we find most worth pursuing. We do not
settle them here; we mark them as the natural sequel to this work.
.
%% rset.colors.tex
%% Red/black colour layer for the two intertwined copies, plus the FM symbols
%% used by rset.mainthm.tex and rset.applications.tex.
%%
%% EVERYTHING routes through \rd (red copy = player) and \bk (black copy =
%% opponent): redefine just those two macros to change the whole scheme (e.g.
%% to make the "black" copy dark blue for legibility).  Every definition is
%% \providecommand-guarded, so it is safe to \input this file from several
%% places and harmless if rset.local already defines these names.
%%
%% Requires \usepackage{xcolor} (already loaded in rset.tex).

% --- the two copies -------------------------------------------------------
\providecommand{\rd}[1]{\textcolor{red}{#1}}     % red copy  (player)
\providecommand{\bk}[1]{\textcolor{black}{#1}}   % black copy (opponent)

% --- theory names and atom supplies --------------------------------------
\providecommand{\rSetX}{\rd{\mathrm{Set}[X]}}    % red   Set[X]
\providecommand{\bSetX}{\bk{\mathrm{Set}[X]}}    % black Set[X]
\providecommand{\rX}{\rd{X}}                      % red   atom supply
\providecommand{\bX}{\bk{X}}                      % black atom supply

% --- empty set ------------------------------------------------------------
\providecommand{\remp}{\rd{\emptyset}}
\providecommand{\bemp}{\bk{\emptyset}}

% --- membership (relations) ----------------------------------------------
\providecommand{\rin}{\mathrel{\rd{\in}}}
\providecommand{\bin}{\mathrel{\bk{\in}}}
\providecommand{\rnin}{\mathrel{\rd{\notin}}}
\providecommand{\bnin}{\mathrel{\bk{\notin}}}

% --- binary operations ----------------------------------------------------
\providecommand{\rcup}{\mathbin{\rd{\cup}}}
\providecommand{\bcup}{\mathbin{\bk{\cup}}}
\providecommand{\rcap}{\mathbin{\rd{\cap}}}
\providecommand{\bcap}{\mathbin{\bk{\cap}}}
\providecommand{\rsm}{\mathbin{\rd{\setminus}}}
\providecommand{\bsm}{\mathbin{\bk{\setminus}}}

% --- comprehension braces of each theory ---------------------------------
\providecommand{\rbr}[1]{\rd{\{}\,#1\,\rd{\}}}
\providecommand{\bbr}[1]{\bk{\{}\,#1\,\bk{\}}}

% --- colour-blind membership (sees through the colour barrier) ------------
\providecommand{\cbin}{\mathrel{\in_{\!\star}}}

% --- FM symbols (shared with rset.stdfmset.tex) --------------------------
\providecommand{\Atoms}{\mathbb{A}}
\providecommand{\Sym}{\mathrm{Sym}}
\providecommand{\pow}{\mathcal{P}}
\providecommand{\supp}{\mathrm{supp}}
\providecommand{\fs}{\mathrm{fs}}
\providecommand{\new}{\mathbin{\reflectbox{\ensuremath{\mathsf{N}}}}}
\providecommand{\atomsof}{\mathrm{atoms}}
\providecommand{\Red}{\rd{\mathsf{R}}}           % the red HF universe
   % red/black colour layer + FM symbols (providecommand-guarded)

\section{Future work: palettes and the generation of colour}
\label{sec:future}

\subsection{Arbitrarily many colours}
\label{sec:future:palette}

Two colours were a convenience, not a necessity. Replace the pair by any set $C$
of \emph{colours}; each $c\in C$ carries its own copy of the theory,
$\mathrm{Set}_c[X_c]$, with its own membership $\in_c$ and its own comprehension
$\{\cdots\}_c$. The only remaining datum is the \emph{nesting discipline}: which
colours may serve as atoms for which. This is a relation on $C$, and it is most
naturally presented as a directed graph --- equivalently a finite-state machine
--- with states $C$ and an edge $c\to c'$ whenever $c'$-sets are admissible
$c$-atoms. The atom supply of each colour is then the disjoint sum over its
outgoing edges,
\[
  X_c \;=\; \bigsqcup_{c\to c'}\ \mathrm{Set}_{c'}[X_{c'}],
\]
and the two-colour theory of this paper is the instance given by the two-state
automaton $\rd{\mathrm{red}} \rightleftarrows \bk{\mathrm{black}}$, with
$\rX=\bSetX$ and $\bX=\rSetX$.

Everything in the paper relativises to this picture. The well-formed cross-colour
nestings are the runs of the automaton; the local/global well-foundedness analysis
of Section~\ref{sec:nwf} reads off its structure, an acyclic colour graph giving a
globally well-founded universe and a cycle giving non-well-foundedness along that
cycle. (A self-loop $c\to c$ would let $c$-sets be their own atoms and collapse
the barrier; the two-cycle $\rd{\mathrm{red}}\rightleftarrows\bk{\mathrm{black}}$
is the smallest discipline that keeps every colour locally well-founded while
making the system globally non-well-founded.) This generalization is routine, and
we record it mainly to fix the slogan: \emph{a reflective set theory is an
automaton on colours}, and the constructions above are its two-state instance.

\subsection{Generating colours dynamically}
\label{sec:future:dynamic}

The genuinely interesting question we leave open. A fixed automaton fixes the
palette in advance. One would instead like colours to be \emph{generated on
demand}: a fresh colour spawned in the course of building a set, the palette
growing as construction proceeds, the automaton unfolding rather than given. The
difficulty is not in admitting infinitely many colours --- a countable palette is
unremarkable --- but in doing so \emph{elegantly}, with a generation discipline
internal to the theory, so that the mechanism producing colours is of a piece with
the mechanism producing sets and names.

There is a suggestive circularity to exploit, which is why we think the question
worth posing rather than dismissing. Generating a fresh colour is generating a
fresh name, in exactly the sense that
Sections~\ref{sec:fm}--\ref{sec:applications} reconstruct: a colour is drawn from
a supply, distinct from those already in play, and otherwise structureless. So a
reflective theory might be made to \emph{generate its own colours} by the very
nominal machinery it models --- the palette drawn, like the atoms, from a stream,
and the colour automaton itself becoming a nominal object in a theory one level
up. Whether this reflexive step can be taken without infinite regress, what its
fixed point looks like, and whether it lands back inside the spectrum of
Section~\ref{sec:nwf}, are the questions we find most worth pursuing. We do not
settle them here; we mark them as the natural sequel to this work.
